% Part: first-order-logic
% Chapter: proof-systems
% Section: introduction

\documentclass[../../../include/open-logic-section]{subfiles}

\begin{document}

\iftag{FOL}
      {\olfileid{fol}{prf}{int}}
      {\olfileid{pl}{prf}{int}}

\olsection{Inleiding}

Een logica heeft meestal zowel een semantiek als
!!a{derivation}ssysteem. De semantiek gaat over begrippen als waarheid,
vervulbaarheid, geldigheid en semantisch gevolg. Een
!!{derivation}ssysteem geeft een methode die alleen naar de vorm van de
symbolen kijkt. Met zo'n methode kunnen we semantisch gevolg en
geldigheid vaststellen. Het systeem is zuiver syntactisch: elke
!!a{derivation} erin is een eindig syntactisch object. Meestal is dat
een rij, of een andere eindige rangschikking, van !!{sentence}s of
!!{formula}s. Bij een goed !!{derivation}ssysteem kan een vaste
procedure van elke gegeven rij of rangschikking van !!{sentence}s of
!!{formula}s controleren of die ``correct'' is.

De eenvoudigste !!{derivation}ssystemen voor de eerste-ordelogica waren
ook de eerste die historisch werden gebruikt. Ze waren
\emph{axiomatisch}. In zo'n systeem geldt een rij !!{formula}s als
!!a{derivation} als elke afzonderlijke !!{formula} aan een van twee
voorwaarden voldoet. Ofwel staat de formule in een vaste verzameling
``axioma's'', ofwel volgt zij met een van een vast aantal
``afleidingsregels'' uit eerdere !!{formula}s in de rij. Een vaste
procedure moet bovendien kunnen controleren of !!a{formula} een axioma
is en of zij volgens een afleidingsregel correct uit andere
!!{formula}s volgt. Axiomatische !!{derivation}ssystemen zijn eenvoudig
te beschrijven en ook eenvoudig te bestuderen als systemen op zichzelf,
dus metatheoretisch. De !!{derivation}s
zelf zijn echter moeilijk te lezen en te begrijpen. Ze zijn ook
moeilijk op te stellen.

Andere !!{derivation}ssystemen zijn bedoeld om !!{derivation}s
gemakkelijker op te stellen of om een voltooide !!{derivation}
gemakkelijker te begrijpen. Voorbeelden zijn natuurlijke deductie,
waarheidsbomen, die ook tableaubewijzen heten, en de sequentencalculus.
Sommige !!{derivation}ssystemen zijn speciaal ontworpen om door een
machine te worden gebruikt. De resolutiemethode is bijvoorbeeld
gemakkelijk in software uit te voeren, al zijn de !!{derivation}s
ervan vrijwel niet te begrijpen. De meeste van deze andere
!!{derivation}ssystemen geven een !!{derivation} weer als een boom van
!!{formula}s en niet als een rij. In zo'n boom kun je gemakkelijker
zien welke delen van de !!a{derivation} afhangen van welke andere
delen.

Voor één logica, bijvoorbeeld de eerste-ordelogica, kunnen de
verschillende !!{derivation}ssystemen dus elk op hun eigen manier
vastleggen wanneer !!a{sentence} een \emph{stelling} is en wanneer
!!a{sentence} uit andere zinnen !!{derivable} is. Dat kan gebeuren met
axiomatische !!{derivation}s, natuurlijke deducties,
sequent!!{derivation}s, waarheidsbomen of resolutieweerleggingen. Welke
methode we ook kiezen, deze relaties moeten passen bij de semantische
begrippen geldigheid en semantisch gevolg. We schrijven $\Proves !A$
voor ``$!A$~is een stelling'' en ``$\Gamma \Proves !A$'' voor
``$!A$~is uit~$\Gamma$ !!{derivable}.'' Hoe we $\Proves$ ook
definiëren, het teken moet hetzelfde verband weergeven als $\Entails$:
\begin{enumerate}
\item $\Proves !A$ dan en slechts dan als $\Entails !A$
\item $\Gamma \Proves !A$ dan en slechts dan als $\Gamma \Entails !A$
\end{enumerate}
De richting ``alleen als'' heet \emph{correctheid}. Een
!!^a{derivation}ssysteem is correct als uit !!{derivability} altijd
semantisch gevolg (of geldigheid) volgt. Elk deugdelijk
!!{derivation}ssysteem moet correct zijn. Aan een incorrect
!!{derivation}ssysteem heb je niets. Het doel van !!a{derivation} is
immers juist een syntactische garantie van geldigheid of semantisch
gevolg te geven. Voor de !!{derivation}ssystemen die we behandelen,
zullen we de correctheid bewijzen.

De omgekeerde richting, ``als'', is ook belangrijk. Deze richting heet
\emph{volledigheid}. Een volledig !!{derivation}ssysteem is sterk
genoeg om te laten zien dat $!A$~een stelling is zodra $!A$~geldig is.
Het kan ook laten zien dat $\Gamma \Proves !A$ zodra
$\Gamma \Entails !A$. Volledigheid is moeilijker te bewijzen. Sommige
logica's hebben geen volledig !!{derivation}ssysteem, maar de
eerste-ordelogica wel. Kurt G\"odel bewees in zijn proefschrift uit
1929 als eerste de volledigheid voor !!a{derivation}ssysteem voor de
eerste-ordelogica.

Bij !!{derivation}ssystemen hoort nog een begrip:
\emph{consistentie}. Een verzameling !!{sentence}s heet inconsistent
als werkelijk alles eruit kan worden !!{derive}d. Anders heet de
verzameling consistent. Inconsistentie is de syntactische tegenhanger
van onvervulbaarheid. Net als onvervulbare verzamelingen zijn
inconsistente verzamelingen !!{sentence}s geen goede theorieën: er is
iets fundamenteel mis mee. Consistente verzamelingen !!{sentence}s
hoeven niet waar of bruikbaar te zijn. Ze halen ten minste wel deze
minimale grens van logische bruikbaarheid. De precieze definitie van
consistentie van verzamelingen !!{sentence}s kan per
!!{derivation}ssysteem verschillen. Net als bij~$\Proves$ willen we
dat consistentie overeenkomt met het semantische begrip
vervulbaarheid. We willen dat altijd geldt: $\Gamma$ is consistent
dan en slechts dan als zij vervulbaar is. De richting ``alleen als''
komt neer op volledigheid: consistentie garandeert vervulbaarheid. De
richting ``als'' komt neer op correctheid: vervulbaarheid garandeert
consistentie. Voor de klassieke eerste-ordelogica zijn deze twee
versies van correctheid en volledigheid inderdaad equivalent.
\end{document}
